\section{Learned Stopping Signals}
\label{sec:problem}

\subsection{Stopping under a search budget}

An episode begins with an information goal $g$ and permits at most $B$
retrieval actions. We index the resulting prefixes by $t=0,\ldots,B$. After
$t$ searches, $E_t$ is the
deduplicated set of retained retrieval units and $b_t=B-t$ is the remaining
budget. We write the environment state as $s_t=(g,E_t,b_t)$. The frozen answer
generator produces $a_t=A(g,E_t)$ at every prefix. A search returns
$o_{t+1}$ and updates the evidence to
$E_{t+1}=\operatorname{dedup}(E_t\cup o_{t+1})$. Because our learned scores
and repetition rule use previous answers, the policy observes the augmented
prefix $x_t=(s_t,a_{0:t})$.

Let $y^\star$ be the reference answer and let $u(a,y^\star)\in[0,1]$ denote
answer quality. We use token F1 to construct the training target and report
exact match (EM) as the primary policy metric on a 0 to 100 scale. Define
$U(x_t)=\mathbb{E}[u(a_t,y^\star)\mid x_t]$. A cost-sensitive policy for this
fixed query-generation and retrieval process ideally solves
\begin{equation}
\max_{\pi}\ \mathbb{E}_{\pi}\!\left[U(x_T)-c_sT\right],
\label{eq:objective}
\end{equation}
where $c_s>0$ is a conceptual search price and $T\le B$ is the stopping time.
We do not estimate $c_s$ or solve this objective. Instead, we select score
thresholds on development trajectories and measure the resulting calls,
tokens, and wall time. The selected threshold is therefore not a monetary
price.

\subsection{Future-best gain}
\label{sec:fbg}

We collect full-budget trajectories with a fixed-budget ReAct logging policy. A
trajectory records the chosen queries, retrieved evidence, and answer at every
prefix:
\[
\xi=(s_0,a_0,o_1,s_1,a_1,\ldots,o_B,s_B,a_B).
\]
For $t<B$, the \emph{future-best gain} is
\begin{equation}
\mathrm{FBG}_t=
\left[\max_{k\in\{t+1,\ldots,B\}}
u(a_k,y^\star)-u(a_t,y^\star)\right]_+,
\qquad \mathrm{FBG}_B=0.
\label{eq:fbg}
\end{equation}
The maximum is available only after the logged future is observed. FBG is thus
a hindsight label for one realized continuation. It is not a Bellman value, a
calibrated value of information, or the gain available under every possible
future query. Clipping also maps a harmful continuation and a neutral one to
the same zero target. Our original cost-adjusted, discounted, and one-step
variants also apply positive-part clipping and remain nonnegative. They do not
test signed supervision. A separate retrospective refit below uses
$\max_{k>t}[u(a_k,y^\star)-u(a_t,y^\star)]\in[-1,1]$, without clipping
(and zero at the terminal prefix). Rerolled continuations probe a different
limitation: unobserved futures outside the logged path.

Our value head estimates
$\fbgpred_t\approx\mathbb{E}[\mathrm{FBG}_t\mid x_t]$. Because the target
looks across all later prefixes, the score can remain positive when the next
retrieval alone does not help. Throughout our paper, \emph{FBG predictor}
means this supervised regressor.

We also define the binary future-improvement label
\begin{equation}
c_t=\mathbf{1}[\mathrm{FBG}_t>0].
\label{eq:future-improvement}
\end{equation}
Our binary future-improvement classifier estimates
$\Pr(c_t=1\mid x_t)$. It predicts whether the logged future contains a
positive gain and does not use an independently annotated stop action.

\subsection{Current correctness and sufficiency}
\label{sec:alternative-signals}

We compare FBG and binary future improvement with two current-state targets.
The error head $\hat r(x_t)$ estimates the probability that $a_t$ is
incorrect. This target asks whether the current answer is usable, whereas the
two future-gain targets ask whether the logged continuation contains a better
answer.

The gold-support sufficiency classifier uses a distinct training label. Let
$z_t=1$ when the lowercased alphanumeric tokens in $E_t$ cover at least 0.8 of
the annotated support-sentence tokens, and let $z_t=0$ otherwise. We train
$\hat p_{\mathrm{suff}}(x_t)$ on $z_t$.
Gold support defines supervision only and is unavailable to the deployed
classifier. We evaluate the nine controller features and a declared
twelve-feature lexical variant. This signal measures evidence coverage, which
is neither answer correctness nor future answer gain.

\subsection{A post-stop selective-risk certificate}
\label{sec:conformal-theory}

The certificate is separate from all stopping policies. After a stopping rule
and $\hat r$ are fixed, a disjoint calibration set of $m$ episodes yields one
first-crossing input and correctness label $(x_{T_i},e_i)$ per episode. Here
$e_i=1$ denotes an incorrect answer. Selection with guaranteed risk
(SGR)~\cite{geifman2017selective} searches nested error-score sets using an
exact binomial upper confidence bound at level
$\delta/\lceil\log_2 m\rceil$. It returns a cutoff $q_\alpha$ only when a
nonempty set satisfies target risk $\alpha$. Otherwise, we set
$q_\alpha=\bot$ and report zero certificate coverage. At each fixed
$\alpha\in\{0.1,0.2,0.3,0.4,0.5,0.6\}$, SGR uses $\delta=0.10$ and the
guarantee below holds separately. They never affect stopping.

\begin{proposition}[Post-stop certificate risk bound]
\label{prop:conformal}
For fixed $\alpha$, assume independent, identically distributed calibration
and test episodes, with the stopping rule and score fixed before calibration
labels. With probability at least $1-\delta$ over calibration, the nonempty
subset returned by SGR
satisfies
\[
\Pr[e=1\mid \hat r(x_T)\le q_\alpha]\le\alpha.
\]
If no candidate meets the corrected bound, no certificate is returned.
\end{proposition}

The guarantee is marginal over episodes rather than per question or subgroup.
The supplement gives the proof sketch and the exact refusal convention.
